\par\Needspace{4\baselineskip}
\begin{center}
\hypertarget{native-060-article043}{}\pdfbookmark[2]{43.}{article-043}43.
\end{center}
Zur weiteren Erläuterung der vorangegangenen Untersuchungen fügen wir ein vollständiges Beispiel für die Berechnung der wahren Anomalie und des Radiusvektors aus der Zeit hinzu. Dazu greifen wir die Zahlen aus Art. 38 wieder auf. Wir setzen also $e=0{,}967\,6456\,7$, $\log q=9{,}765\,6500$, $t=63{,}54400$ und leiten daraus zunächst die Konstanten $\log\alpha=0{,}305\,2357$, $\log\beta=8{,}221\,7364$, $\log\gamma=0{,}002\,8755$ ab.

\SourceLeaf{061}Daraus ergibt sich $\log\alpha t=2{,}108\,3102$, wozu die \textsc{Barker}sche Tafel den Näherungswert $w=99^0\,6^{\prime}$ liefert. Daraus folgt $A=0{,}02292$, und unsere Tafel liefert $\log B=0{,}000\,0040$. Das korrigierte Argument, mit dem man in die \textsc{Barker}sche Tafel eingeht, ist damit $=\log\dfrac{\alpha t}{B}=2{,}108\,3062$. Ihm entspricht $w=99^0\,6^{\prime}\,13^{\prime\prime}{,}18$. Die weitere Rechnung verläuft wie folgt.
\begin{center}
\begin{tabular}{@{}ll@{}}
\makebox[18em][l]{$\log\operatorname{tang}\tfrac{1}{2}w^2$\hfill .\,.\,.\,.\,.\,.\,.\,.\,.} & $0{,}138\,5934$\\[.3em]
\makebox[18em][l]{$\log\beta$\hfill .\,.\,.\,.\,.\,.\,.\,.\,.\,.\,.\,.\,.\,.} & $8{,}221\,7364$\\[.3em]
\hline
\makebox[18em][l]{$\log A$\hfill .\,.\,.\,.\,.\,.\,.\,.\,.\,.\,.\,.\,.} & $8{,}360\,3298$\\[.3em]
$A$\hfill $=$ & $0{,}022\,9260\,8$\\[.3em]
\multicolumn{2}{@{}l@{}}{Daraus ergibt sich derselbe $\log B$ wie zuvor}\\[.3em]
$C$\hfill $=$ & $0{,}000\,0242$\\[.3em]
$1-\tfrac{4}{5}A+C$\hfill $=$ & $0{,}981\,6833$\\[.3em]
$1+\tfrac{1}{5}A+C$\hfill $=$ & $1{,}004\,6094$
\end{tabular}
\end{center}
\begin{center}
\begin{tabular}{@{}ll@{}}
\makebox[20em][l]{$\log\operatorname{tang}\tfrac{1}{2}w$\hfill .\,.\,.\,.\,.\,.\,.\,.\,.\,.} & $0{,}069\,2967$\\[.3em]
\makebox[20em][l]{$\log\gamma$\hfill .\,.\,.\,.\,.\,.\,.\,.\,.\,.\,.\,.\,.\,.\,.} & $0{,}002\,8755$\\[.3em]
\makebox[20em][l]{$\tfrac{1}{2}\operatorname{Comp.}\log(1-\tfrac{4}{5}A+C)$\hfill .} & $0{,}004\,0143$\\[.3em]
\hline
\makebox[20em][l]{$\log\operatorname{tang}\tfrac{1}{2}v$\hfill .\,.\,.\,.\,.\,.\,.\,.\,.\,.\,.} & $0{,}076\,1865$\\[.3em]
$\tfrac{1}{2}v$\hfill $=$ & $50^0\quad0^{\prime}\quad0^{\prime\prime}$\\[.3em]
$v$\hfill $=$ & $100\quad0\quad0$\\[.3em]
\makebox[20em][l]{$\log q$\hfill .\,.\,.\,.\,.\,.\,.\,.\,.\,.\,.\,.\,.\,.\,.} & $9{,}765\,6500$\\[.3em]
\makebox[20em][l]{$2\operatorname{Comp.}\log\cos\tfrac{1}{2}v$\hfill .\,.\,.\,.\,.\,.} & $0{,}383\,8650$\\[.3em]
\makebox[20em][l]{$\log(1-\tfrac{4}{5}A+C)$\hfill .\,.\,.\,.\,.\,.\,.} & $9{,}991\,9714$\\[.3em]
\makebox[20em][l]{$\mathrm{C.}\log(1+\tfrac{1}{5}A+C)$\hfill .\,.\,.\,.\,.} & $9{,}998\,0028$\\[.3em]
\hline
\makebox[20em][l]{$\log r$\hfill .\,.\,.\,.\,.\,.\,.\,.\,.\,.\,.\,.\,.\,.\,.} & $0{,}139\,4892$
\end{tabular}
\end{center}
\noindent Hätte man in dieser Rechnung den Faktor $B$ ganz vernachlässigt, so wäre die wahre Anomalie um nur $1^{\prime\prime}$ zu groß ausgefallen.
